% =============================================================================
% components/experience.tex — experiência profissional
% -----------------------------------------------------------------------------
%   \cvJob{<cargo>}{<organização>}{<período>}
%     \cvItem{...}
%     \cvItem{...}
%
% Gera:  **Cargo** - Organização          (Noto Serif Bengali 11)
%        (Período)                        (Calibri 10, #666666)
%
% Opção (em main.tex):
%   \cvJobDateInlinefalse   período numa linha própria, abaixo do cargo (padrão)
%   \cvJobDateInlinetrue    período à direita, na 1ª linha do cargo
%                           (só recomendado se todos os cargos couberem em 1 linha)
% =============================================================================

\newif\ifcvJobDateInline
\cvJobDateInlinefalse

% distância mínima entre o texto do cargo e o período
\newlength\cvJobDateGap  \setlength\cvJobDateGap{12bp}

\cvDefineStyle{job-title}{
  metrics = serif, font = \cvSerif, size = 11, line-height = 1.3,
  before = 12bp, after = 2bp, align = justify}

\cvDefineStyle{job-period}{
  size = 10, line-height = 1.0, after = 6bp, color = cvMuted}

\newsavebox\cvJobDateBox

\newcommand\cvJob[3]{%
  \ifcvJobDateInline
    \sbox\cvJobDateBox{\cvSans\cvFontSize{10}\color{cvMuted}(#3)}%
    \cvParagraph[style=job-title, after=6bp,
                 right-indent=\dimexpr\wd\cvJobDateBox+\cvJobDateGap\relax]{%
      \rlap{\hskip\dimexpr\linewidth-\wd\cvJobDateBox\relax
            \usebox\cvJobDateBox}%
      \textbf{#1} - #2}%
  \else
    \cvParagraph[style=job-title]{\textbf{#1} - #2}%
    \cvParagraph[style=job-period]{(#3)}%
  \fi}
